% Options for packages loaded elsewhere
\PassOptionsToPackage{unicode}{hyperref}
\PassOptionsToPackage{hyphens}{url}
\PassOptionsToPackage{dvipsnames,svgnames,x11names}{xcolor}
%
\documentclass[
  11pt,
  a4paper,
]{article}
\usepackage{amsmath,amssymb}
\usepackage{iftex}
\ifPDFTeX
  \usepackage[T1]{fontenc}
  \usepackage[utf8]{inputenc}
  \usepackage{textcomp} % provide euro and other symbols
\else % if luatex or xetex
  \usepackage{unicode-math} % this also loads fontspec
  \defaultfontfeatures{Scale=MatchLowercase}
  \defaultfontfeatures[\rmfamily]{Ligatures=TeX,Scale=1}
\fi
\usepackage{lmodern}
\ifPDFTeX\else
  % xetex/luatex font selection
  \setmainfont[]{Latin Modern Roman}
  \setmathfont[]{Latin Modern Math}
\fi
% Use upquote if available, for straight quotes in verbatim environments
\IfFileExists{upquote.sty}{\usepackage{upquote}}{}
\IfFileExists{microtype.sty}{% use microtype if available
  \usepackage[]{microtype}
  \UseMicrotypeSet[protrusion]{basicmath} % disable protrusion for tt fonts
}{}
\makeatletter
\@ifundefined{KOMAClassName}{% if non-KOMA class
  \IfFileExists{parskip.sty}{%
    \usepackage{parskip}
  }{% else
    \setlength{\parindent}{0pt}
    \setlength{\parskip}{6pt plus 2pt minus 1pt}}
}{% if KOMA class
  \KOMAoptions{parskip=half}}
\makeatother
\usepackage{xcolor}
\usepackage[top=25mm,bottom=24mm,left=27mm,right=27mm]{geometry}
\setlength{\emergencystretch}{3em} % prevent overfull lines
\providecommand{\tightlist}{%
  \setlength{\itemsep}{0pt}\setlength{\parskip}{0pt}}
\setcounter{secnumdepth}{-\maxdimen} % remove section numbering
% Mathematics-focused preprint layout: monochrome, Latin Modern OpenType text and math.
% Latin Modern fonts are supplied by a standard TeX installation, not embedded in this source archive.
\usepackage{microtype}
\usepackage{setspace}
\usepackage{titlesec}
\usepackage{fancyhdr}
\usepackage{needspace}
\setstretch{1.13}
\setlength{\parindent}{1.15em}
\setlength{\parskip}{3.25pt plus 1.0pt minus .6pt}
\titlespacing*{\section}{0pt}{17pt plus 4pt minus 2pt}{7pt plus 2pt}
\titlespacing*{\subsection}{0pt}{13pt plus 3pt minus 1pt}{5pt plus 1pt}
\titlespacing*{\subsubsection}{0pt}{10pt plus 2pt minus 1pt}{3pt plus 1pt}
\titleformat{\section}{\normalfont\bfseries\large}{\thesection}{0.7em}{}
\titleformat{\subsection}{\normalfont\bfseries\normalsize}{\thesubsection}{0.6em}{}
\titleformat{\subsubsection}{\normalfont\bfseries\normalsize}{\thesubsubsection}{0.6em}{}
\AtBeginDocument{%
  \setlength{\abovedisplayskip}{12pt plus 3pt minus 2pt}%
  \setlength{\belowdisplayskip}{12pt plus 3pt minus 2pt}%
  \setlength{\abovedisplayshortskip}{9pt plus 2pt minus 1pt}%
  \setlength{\belowdisplayshortskip}{10pt plus 2pt minus 1pt}%
  \pagestyle{plain}%
}
\setlength{\emergencystretch}{4em}
\setcounter{secnumdepth}{0}
\ifLuaTeX
  \usepackage{selnolig}  % disable illegal ligatures
\fi
\usepackage{bookmark}
\IfFileExists{xurl.sty}{\usepackage{xurl}}{} % add URL line breaks if available
\urlstyle{same}

\usepackage{longtable,booktabs,array,calc}
\usepackage{fvextra}
\RecustomVerbatimEnvironment{verbatim}{Verbatim}{breaklines=true,breakanywhere=true,fontsize=\footnotesize}
\usepackage{upquote}

\hypersetup{
  pdftitle={Division-polynomial covariance and Kummer sequences on secp256k1},
  pdfauthor={Louis Desbiens (human initiator and curator)},
  pdfsubject={AI-developed mathematical research artifact; independently unverified candidate results},
  colorlinks=true,
  linkcolor={black},
  filecolor={black},
  citecolor={black},
  urlcolor={black},
  pdfcreator={LaTeX via pandoc}}

\title{Division-polynomial covariance and Kummer sequences on secp256k1}
\usepackage{etoolbox}
\makeatletter
\providecommand{\subtitle}[1]{%
  \apptocmd{\@title}{\par {\large #1 \par}}{}{}
}
\makeatother
\subtitle{Exact character targets, disjoint branch layers, and limits of finite-orbit inference}
\author{}
\date{Public mathematical research artifact \textbar{} 31 October 2026}
\begin{document}
\maketitle

\textbf{Human initiator and curator:} Louis Desbiens\\
\textbf{AI mathematical research contributor:} OpenAI's ChatGPT\\
\textbf{Document classification:} Public AI-developed mathematical
research artifact -- independently unverified candidate results\\
\textbf{Institutional affiliation:} None claimed\\
\textbf{Manuscript license:} Creative Commons Attribution 4.0 International (CC BY 4.0)

\textbf{Status and reading guide.} This document records a
human-initiated, substantially AI-developed mathematical investigation.
The human curator directed the objectives but does not independently
certify its advanced mathematical arguments. The novel claims and
supplied proofs are candidate results awaiting independent specialist
verification; public availability does not certify correctness,
completeness or originality. This artifact proposes Kummer covariance
obstructions, recurrence and word-occurrence bounds, and higher-genus
quotient/isogeny constructions. The integer threshold checks do not
verify the underlying character-sum estimates, geometric independence,
conductor assumptions or isogeny/descent arguments. Those steps require
independent mathematical examination. No unrestricted
private-key-recovery claim is made.

\section{Abstract}\label{abstract}

The division-polynomial investigation is organized around an exact
covariance equation: a base-field sextic character of a rational
function is sought that reproduces the sextic character of a hidden
scalar. On secp256k1, a primitive degree-three CM endomorphism turns the
question into a cyclic-character problem. The manuscript proposes a
divisor calculation for a division-polynomial covariance ratio,
identifies a cubic atom, and investigates an obstruction to the
separable single-division-polynomial family as exact scalar-character
extractors. A proposed conductor argument yields large effective divisor
support for products. Along the primitive orbit, the cubic atom defines
a ternary sequence with nontrivial holonomy; proposed disjoint branch
layers of sizes \(18\cdot3^i\) lead to recurrence and block-occurrence
estimates. Subject to the stated geometric and character-sum hypotheses,
the manuscript proposes exclusion of nonzero affine recurrences of
length at most 78 and occurrence of every ternary word of length at most
39. It also develops a cubic-extension Kummer-symbol model, genus-seven
covering geometry, candidate quotient curves, a Jacobian decomposition,
and a degree-three isogeny between elliptic quotients of a reciprocal
genus-two model. Independence and descent assumptions are separated
explicitly. These proposals concern restricted observables, not general
private-key recovery, and require specialist examination.

\textbf{Keywords:} elliptic divisibility sequence; division polynomial;
Kummer cover; cubic character; CM endomorphism; cyclic sequence;
secp256k1.

\section{1. Setting and conventions}\label{setting-and-conventions}

The relationship between elliptic divisibility sequences and discrete
logarithms is established prior work {[}4{]}. Here the question is
narrower: can a rational function's sextic character reproduce the
multiplicative sextic character of an unknown scalar, at every nonzero
point of one fixed prime-order group? The exploratory investigation used
numerical evaluations to test the initial intuition; the exact
covariance equation below replaces any statistical criterion in the
mathematical argument.

Use the SEC 2 curve {[}5{]}

\[
E/\mathbf F_p:y^2=x^3+7,\qquad p=2^{256}-2^{32}-977.
\]

The group order \(n\) and the two oriented cube roots are specified by
concatenating each pair of hexadecimal lines:

\begin{verbatim}
n      = FFFFFFFFFFFFFFFFFFFFFFFFFFFFFFFE
         BAAEDCE6AF48A03BBFD25E8CD0364141
beta   = 7ae96a2b657c07106e64479eac3434e9
         9cf0497512f58995c1396c28719501ee
lambda = 5363ad4cc05c30e0a5261c028812645a
         122e22ea20816678df02967c1b23bd72
\end{verbatim}

SEC 2 specifies a generator \(G\) of this prime-order group. The
identity \(\phi(G)=[\lambda]G\), where \(\phi(x,y)=(\beta x,y)\), is
checked by the ordinary group law on its published coordinates. Hence
\(\phi=[\lambda]\) on all rational points. We fix this orientation,
rather than interchanging the two cube roots implicitly.

Put \(H=\{\pm1,\pm\lambda,\pm\lambda^2\}\), \(h=2+\lambda\),
\(\Psi=2+\phi=1-\phi^2\), and \(M=(n-1)/6\). Appendix A prints the
complete certificate that \(h\) has order \(n-1\). The map \(\Psi\) is
separable of degree three: its kernel is
\(\{O,(0,\sqrt7),(0,-\sqrt7)\}\), and its differential is nonzero.
Alternatively its Eisenstein norm is three. The chord from \(P\) to
\(-\phi^2P\) gives

\[
x(\Psi P)=-\frac{\beta}{3}\frac{x^3+28}{x^2},\qquad
F(t)=-\frac{(t+28)^3}{27t^2},\quad t=x^3. \qquad (1.1)
\]

Indeed the slope is \(-2y/[x(\beta^2-1)]\), and
\((\beta^2-1)^2=-3\beta^2\). Cubing the abscissa gives
\(t(\Psi P)=F(t(P))\). The identity \(\Psi^2=-[3]\phi^2\) follows by
expanding \((2+\phi)^2\) and using \(1+\phi+\phi^2=0\).

The coordinate \(t\) is a complete sixfold quotient: equality of two
cubes of abscissae gives \(x(Q)=\beta^j x(P)\), and the curve equation
then gives \(y(Q)=\pm y(P)\). No nonidentity rational point has a
nontrivial stabilizer, since that would give two- or three-torsion.
Primitivity of \(h\) therefore makes \(\Psi\) a single \((n-1)\)-cycle
on nonzero points and \(F\) a single \(M\)-cycle on their quotient.

For nonzero field elements set

\[
\eta(z)=z^{(p-1)/6},\qquad \chi_3(z)=z^{(p-1)/3}=\eta(z)^2.
\]

The values here are initially roots of unity in \(\mathbf F_p\).
Whenever absolute values or character sums are used, identify \(\beta\)
with \(e^{2\pi i/3}\) and \(-1\) with the usual complex number \(-1\).
This specifies a compatible complex character. Let
\(\iota:\mu_6(\mathbf F_n)\to\mu_6(\mathbf F_p)\) send \(\lambda\) to
\(\beta\) and \(-1\) to \(-1\). Appendix A verifies

\[
h^M=-\lambda^2,\quad M\equiv1\pmod3,\quad p\equiv7\pmod9,\quad
\chi_3(\beta/81)=1. \qquad (1.2)
\]

All multiplication indices below are positive integers prime to \(p\);
indices divisible by \(n\) are also excluded when evaluating a division
polynomial on every nonidentity rational point.

\section{2. The exact sextic target and
covariance}\label{the-exact-sextic-target-and-covariance}

Let \(P_d=[d]G\) for \(d\in\mathbf F_n^*\) and define the scalar label

\[
C(P_d)=\iota(d^M).
\]

The proposed exact extractor is a rational function
\(g\in\mathbf F_p(E)^*\), finite and nonzero at every \(P_d\), with

\[
\eta(g(P_d))=C(P_d). \qquad (2.1)
\]

Unlike its covariance ratio, this label is not invariant under all
sixfold scalar symmetries. With \(\xi=\iota(h^M)=-\beta^2\), it
satisfies \(C(\Psi P)=\xi C(P)\).

\textbf{Proposition 2.1 (orbit covariance).} Put \(R_g=(g\circ\Psi)/g\).
A necessary condition for (2.1) is

\[
\eta(R_g(P))=\xi. \qquad (2.2)
\]

Equation (2.2) is required at every nonidentity rational point.
Conversely, (2.2) together with (2.1) at one point implies (2.1)
everywhere on that set.

\textbf{Proof.} Taking the ratio of (2.1) at \(\Psi P\) and \(P\) gives
(2.2). Conversely, propagation along the single \((n-1)\)-cycle of
\(\Psi\) determines the label from one anchor. \(\square\)

Thus the exact target has a concrete mathematical meaning. A weak
correlation on sampled points, or an identity valid only after knowing
\(d\), is not a substitute for (2.2).

\section{3. Division-polynomial covariance
divisors}\label{division-polynomial-covariance-divisors}

Let \(\psi_m\) be the standard division polynomial. Its divisor, for
\(p\nmid m\), is {[}1{]}

\[
\operatorname{div}(\psi_m)=[m]^*(O)-m^2(O). \qquad (3.1)
\]

Over an algebraic closure, write \(\tau=(0,\sqrt7)\). The kernel of
\(\Psi\) is \(\{O,\tau,-\tau\}\), and \(\tau\) has order three. Define

\[
R_m(P)=\frac{\psi_m(\Psi P)}{\psi_m(P)}.
\]

\textbf{Proposition 3.1.} The covariance divisor is

\[
\operatorname{div}(R_m)=[m]^*((\tau)+(-\tau))-m^2((\tau)+(-\tau)). \qquad (3.2)
\]

Moreover \(R_m\) is invariant under negation and \(\phi\), so it belongs
to \(\mathbf F_p(t)\).

\textbf{Proof.} Pull (3.1) back by \(\Psi\) and subtract (3.1). Since
the two endomorphisms commute and \(\Psi^*(O)=(O)+(\tau)+(-\tau)\), the
terms involving \((O)\) cancel and give (3.2). The division-polynomial
transformation factors under negation and the cubic automorphism are
constants; those constants cancel in the ratio. Thus \(R_m\) descends
through the sixfold quotient. \(\square\)

If \(n\nmid m\), neither \(P\) nor \(\Psi P\) is killed by \([m]\) for a
nonidentity rational \(P\). Hence \(R_m\) is finite and nonzero at every
such point. The standard composition identity gives the useful exact
relation

\[
R_{ab}(P)=R_a([b]P)R_b(P)^{a^2}. \qquad (3.3)
\]

This formula reduces character behavior of composite indices to atomic
factors evaluated at known multiples. It does not reveal the unknown
scalar of \(P\).

\section{4. The cubic atom and the single-factor
obstruction}\label{the-cubic-atom-and-the-single-factor-obstruction}

Since \(\psi_3=3x(t+28)\), substitution of (1.1) gives

\[
R_3(t)=\frac{\beta}{81}\frac{f(t)}{t^3},\qquad
f(t)=t^3-672t^2+2352t+21952. \qquad (4.1)
\]

To verify (4.1), substitute (1.1) into \(3x(t+28)\) and divide by
\(3x(t+28)\); the numerator factorization is \((t+28)^3-756t^2=f(t)\),
and the remaining factors give the displayed constant and \(t^3\)
denominator. By (1.2),

\[
A(P):=\chi_3(R_3(P))=\chi_3(f(t(P))). \qquad (4.2)
\]

\textbf{Proposition 4.1.} No function \(g=c\psi_m\), with
\(c\in\mathbf F_p^*\), \(p\nmid m\), and \(n\nmid m\), satisfies the
exact extractor equation (2.1).

\textbf{Proof.} First impose the symmetries of the target. Since \(M\)
is even, \(C(-P)=C(P)\). The identities
\(\psi_m(-P)=(-1)^{m-1}\psi_m(P)\) and \(\eta(-1)=-1\) force \(m\) to be
odd. Also \(C(\phi P)=\beta C(P)\). The standard weight transformation
of division polynomials is \(\psi_m(\phi P)=\beta^{2(m^2-1)}\psi_m(P)\);
applying \(\eta\) forces \(3\mid m\). Explicitly,
\((p-1)/6\equiv1\pmod3\), so \(\eta(\beta)=\beta\); the exponent
\(2(m^2-1)\) is zero modulo three when \(3\nmid m\) and one when
\(3\mid m\). Thus \(m=3b\) with \(m\equiv3\pmod6\).

Equation (3.3) implies

\[
\chi_3(R_{3b}(P))=\chi_3(R_3([b]P)), \qquad (4.3)
\]

because the remaining exponent is nine. Multiplication by \(b\) permutes
the nonidentity rational group. Therefore (2.2) would force \(A\) to be
constant there, with value \(\xi^2=\beta\).

The divisor of \(R_3\) has eighteen simple zeros and poles of order nine
at \(\tau\) and \(-\tau\). Its cubic Kummer sheaf is geometrically
nontrivial, ramified at exactly eighteen geometric points. The tame
character-sum bound on a genus-one curve gives

\[
\left|\sum_{P\in E(\mathbf F_p)}\chi_3(R_3(P))\right|\leq18\sqrt p. \qquad (4.4)
\]

Here the function is interpreted by its finite nonzero value at \(O\);
all geometric branch points are nonrational. If the \(n-1\) nonidentity
terms were constant, the full sum would have magnitude at least \(n-2\).
But \(n-2>18\sqrt p\). This contradiction proves the proposition.
\(\square\)

\textbf{Lemma 4.2 (the character-sum input).} Let \(R\) define a
finite-order rank-one Kummer local system on a smooth genus-one curve
over \(\mathbf F_p\), with exactly \(s>0\) geometric ramification
points, all tame and with nontrivial inertia. Give its unramified
extension the usual Frobenius trace at every remaining rational point.
Then the sum of its traces has absolute value at most \(s\sqrt p\).

\textbf{Proof.} Remove precisely the ramification points and denote the
resulting open curve by \(U\). Nontrivial geometric ramification
excludes a constant constituent, so
\(H_c^0(U,\mathcal L)=H_c^2(U,\mathcal L)=0\). The tame
Euler-characteristic formula gives \(\chi_c(U,\mathcal L)=2-2g-s=-s\),
hence \(\dim H_c^1=s\). The finite-order local system is pure of weight
zero; the curve trace formula and the Weil weight bound give eigenvalues
of absolute value at most \(\sqrt p\) in \(H_c^1\) {[}2,3{]}. As inertia
has no invariants at the removed points, its middle-extension traces
there are zero. Taking the trace proves the assertion. At a zero or pole
whose valuation is divisible by the character order, use the unramified
local extension, not an artificial zero value. \(\square\)

For (4.4), all rational points, including \(O\), are unramified and
\(R_3\) is a local unit there. Removing \(O\) changes a sum by at most
one. This accounts explicitly for the constants \(n-2\) and \(+1\) used
later.

The result closes this separable single-factor family even for enormous
\(m\). It does not close every product of factors or every compact
arithmetic circuit.

\section{5. Products and effective
support}\label{products-and-effective-support}

Consider

\[
g=c\prod_{i=1}^s\psi_{m_i}^{e_i},
\]

where \(p\nmid m_i\), \(n\nmid m_i\), and exponents may be reduced
modulo six for the character question. The same symmetry tests give the
necessary congruences

\[
\sum_{2\mid m_i}e_i\equiv0\pmod2,\qquad
\sum_{3\mid m_i}e_i\equiv1\pmod3. \qquad (5.1)
\]

Define \(s_6(R_g)\) as the number of geometric points where the
valuation of \(R_g\) is not divisible by six. This is the effective
branch support for its sextic Kummer character, rather than its
unreduced polynomial degree.

\textbf{Proposed Theorem 5.1 (support obstruction).} Any exact product
extractor satisfying (2.1) must satisfy

\[
n-2\leq s_6(R_g)\sqrt p,\qquad
s_6(R_g)\leq2\sum_{e_i\not\equiv0\,(6)}(m_i^2+1). \qquad (5.2)
\]

\textbf{Proof.} The second bound follows by taking the union of the two
degree-\(m_i^2\) pullback supports and the two possible poles in (3.2);
cancellations only decrease the support. For the first bound, a
geometrically nontrivial rank-one Kummer sheaf gives the same genus-one
estimate as in Section 4, with \(s_6\) in place of eighteen.

One must justify nontriviality instead of assuming it. Suppose every
valuation of \(R_g\) were divisible by six. The finite divisor
\(D=\operatorname{div}(g)\), taken modulo six, would satisfy
\(\Psi^*D=D\). Since \(\Psi\) is etale of degree three, the support of
the pullback has three times as many geometric points as that of \(D\).
Thus this finite support must be empty. Write
\(\operatorname{div}(g)=6D_0\) with \(D_0\) rational. Its class is
rational six-torsion in \(\operatorname{Pic}^0(E)=E\), but the rational
group has prime order coprime to six. Consequently \(D_0\) is principal.
A function with this rational divisor can be chosen over
\(\mathbf F_p\): its one-dimensional space of sections is
Frobenius-stable, and a scalar rescaling removes the Frobenius
multiplier (equivalently, use Hilbert 90 over a finite extension). Thus
\(g=c_0g_0^6\) with \(c_0\in\mathbf F_p^*\). Its sextic character is
constant, contradicting (2.1). Hence \(s_6>0\) for an exact extractor,
and the Kummer estimate applies. Covariance makes all nonidentity
summands equal; adding the identity term gives the lower bound \(n-2\).
\(\square\)

The bound is on effective divisor complexity. Large indices can have
short evaluation recurrences, so (5.2) is not a general running-time
lower bound.

\section{6. A cubic splitting field and geometric
independence}\label{a-cubic-splitting-field-and-geometric-independence}

The polynomial in (4.1) has

\[
\operatorname{disc}(f)=5334336^2=2^{12}3^{10}7^6,
\qquad f(-7)=-27783. \qquad (6.1)
\]

To prove irreducibility explicitly, let \(K=\mathbf F_{p^3}\) and choose
\(\rho\) with \(\rho^3=19+\beta\). The exact residue identity
\((19+\beta)^{(p-1)/3}=\beta\) gives \(\rho^p=\beta\rho\). One root of
\(f\) is

\[
\alpha=224+84\left(\rho+\frac7\rho\right),\qquad
f(t)=N_{K/\mathbf F_p}(t-\alpha). \qquad (6.2)
\]

\textbf{Verification of the root and its degree.} The identity
\((19+\beta)(19+\beta^2)=343\) gives \(\rho^3+(7/\rho)^3=37\). Thus
\(u=\rho+7/\rho\) satisfies \(u^3-21u-37=0\), and direct substitution
yields

\[
f(224+84u)=84^3(u^3-21u-37).
\]

The residue check in Appendix A makes \(19+\beta\) a noncube, so
\(1,\rho,\rho^2\) is a basis of \(K\) over \(\mathbf F_p\). If
\(\rho+7/\rho\) were Frobenius-fixed, then
\((\beta-1)\rho+7(\beta^2-1)/\rho=0\), forcing \(\rho^2\in\mathbf F_p\),
a contradiction. Hence \(\alpha\) has degree three. Its three Frobenius
conjugates are exactly the roots of the monic cubic \(f\), proving both
irreducibility and (6.2).

Consequently, with \(\widetilde\chi_3(z)=z^{(p^3-1)/3}\),

\[
A(P)=\widetilde\chi_3(t(P)-\alpha). \qquad (6.3)
\]

This is a linear Kummer symbol over a small extension. It is not itself
the usual three-isogeny descent function. Over a field containing
\(s=\sqrt7\), that descent-type function is \(y-s\), with divisor
\(3(\tau)-3(O)\). In contrast, \(t-\alpha\) has six simple zeros.

\textbf{Proposition 6.1.} The smooth normalization of

\[
\mathcal C:\quad y^2=x^3+7,\qquad z^3=x^3-\alpha
\]

has genus seven. Its pullback by \([b]\), for \(b>0\) and \(p\nmid b\),
has genus \(6b^2+1\).

\textbf{Proof.} The cubic Kummer cover of \(E\) has six simple branch
points and an unramified pole of order six at \(O\). Riemann--Hurwitz
gives \(2g-2=6(3-1)=12\). An etale pullback by \([b]\) gives \(6b^2\)
branch points, yielding \(g=6b^2+1\). \(\square\)

\textbf{Proposition 6.2 (positive-index independence).} For distinct
positive integers \(b_1,\ldots,b_s\), all prime to \(p\), the classes of

\[
h_b(P)=t([b]P)-\alpha
\]

are linearly independent in \(\bar K(E)^*/\bar K(E)^{*3}\). Their cube
roots define a cover with group \((\mathbf Z/3)^s\).

\textbf{Proof.} The six zeros of \(t-\alpha\) are points of exact order
nine: they belong to \(\ker\Psi^3\setminus\ker\Psi^2\). For the largest
index \(b\), choose a preimage \(Q\) of one such zero under \([b]\) with
exact order \(9b\). Such a preimage exists by choosing compatible bases
of the prime-to-\(p\) Tate modules: a primitive order-nine vector lifts
under \([b]\) to an order-\(9b\) vector, and the prime-to-three factors
of \(b\) can be supplied in the other torsion components. If \([c]Q\)
were a zero for a smaller positive index \(c\), its order would be nine,
forcing \(\gcd(c,9b)=b\) and hence \(b\mid c\), impossible. Nor is it a
pole for a smaller index. The valuation at \(Q\) isolates the exponent
of \(h_b\) in any proposed cube relation. It must be zero modulo three.
Descending through the indices proves independence. \(\square\)

The restriction to positive indices is explicit: \(h_b=h_{-b}\), so
unrestricted ``distinct indices'' would be false. Also, functions with
distinct geometric indices may agree on the finite rational group when
the indices coincide modulo \(n\). Geometric Kummer independence must
never be confused with independence of finite samples or with the
absence of every finite-orbit character identity.

\section{7. The ternary sequence and its
holonomy}\label{the-ternary-sequence-and-its-holonomy}

Set

\[
P_r=\Psi^rG=[h^r]G,\quad t_r=F^r(t(G)),\quad
A(P_r)=\beta^{a_r},\qquad a_r\in\mathbf F_3.
\]

Since \(A\) is sixfold invariant, \((a_r)\) is \(M\)-periodic. If
\(b=h^s\pmod n\), then \(A([b]P_r)=A(P_{r+s})\). Thus scalar pullbacks
of this atom are shifts of a single cyclic sequence, and a
multiplicative character relation becomes an affine relation among its
shifts.

Define \(\chi_3(\psi_3(P_r))=\beta^{c_r}\). The ratio defining \(R_3\)
gives

\[
a_r=c_{r+1}-c_r. \qquad (7.1)
\]

The \(c_r\) are not \(M\)-periodic. Since \(h^M=-\lambda^2\) and
\(\psi_3(\phi P)=\beta\psi_3(P)\), the orientation identities yield

\[
c_{r+M}=c_r+1,\qquad \sum_{r=0}^{M-1}a_r=1\pmod3. \qquad (7.2)
\]

For the sign and exponent in (7.2), \(P_{r+M}=-\phi^2P_r\). The
polynomial \(\psi_3\) is even under negation and has \(\phi\)-weight
\(\beta\), so \(\psi_3(P_{r+M})=\beta^2\psi_3(P_r)\). Since
\((p-1)/3\equiv2\pmod3\), applying \(\chi_3\) multiplies by
\(\beta^{4}=\beta\). This gives the increment one. Summing (7.1) over a
full quotient cycle telescopes to that increment, rather than to zero.

With \(d_r=c_r-r\) and \(b_r=a_r-1\), one obtains genuinely periodic
sequences satisfying

\[
b_r=d_{r+1}-d_r,\qquad \sum_rb_r=0. \qquad (7.3)
\]

This is nontrivial cyclic holonomy, not a globally periodic primitive
\(c_r\).

For clarity, fix the backward cyclic shift \((Sb)_r=b_{r-1}\) and write
\(B(X)=\sum_rb_rX^r\) in \(\mathbf F_3[X]/(X^M-1)\). Its minimal
annihilator is

\[
m_b(X)=\frac{X^M-1}{\gcd(B(X),X^M-1)}. \qquad (7.4)
\]

Indeed, \(E(S)b=0\) is equivalent to \((X^M-1)\mid E(X)B(X)\). A
forward-shift convention replaces the annihilator by its reciprocal; the
convention is therefore part of the statement. Since \(3\nmid M\),
\(X^M-1\) is square-free, and \(B(1)=0\).

\textbf{Open problem.} Does \(\gcd(B(X),X^M-1)=X-1\) hold? This
maximal-complexity possibility is not proved here. The local results
below do not settle that global question.

\section{8. Disjoint branch layers and recurrence
bounds}\label{disjoint-branch-layers-and-recurrence-bounds}

Let \(K_j=\ker\Psi^j\) over an algebraic closure, so \(|K_j|=3^j\). The
cubic Kummer branch support of \(R_3\) is

\[
B_0=K_3\setminus K_2.
\]

For \(r_i(P)=R_3(\Psi^iP)\), the effective cubic support is

\[
B_i=K_{i+3}\setminus K_{i+2},\qquad |B_i|=18\cdot3^i. \qquad (8.1)
\]

To justify (8.1), observe that \(K_2=E[3]\) by \(\Psi^2=-[3]\phi^2\). In
(3.2) with \(m=3\), the eighteen preimages of \(\tau\) and \(-\tau\) are
precisely \(K_3\setminus K_2\), with multiplicity one; the only poles
are \(\tau,-\tau\), of order nine. Pullback by the etale map \(\Psi^i\)
gives \(\Psi^{-i}(K_j)=K_{i+j}\) and preserves the local multiplicities.
Every pulled-back pole still has valuation divisible by three.

These supports are pairwise disjoint. The poles have valuations
divisible by three and do not enter the cubic conductor. Therefore a
nonzero exponent vector \(e_i\in\{0,1,2\}\) gives exactly

\[
s\left(\prod_i r_i^{e_i}\right)=18\sum_{e_i\neq0}3^i. \qquad (8.2)
\]

There is no cancellation between different levels modulo cubes.

\textbf{Proposed Theorem 8.1 (short affine recurrences).} A nonzero
relation

\[
\sum_{i=0}^{k-1}e_i a_{r+i}=c\qquad\text{for every }r
\]

requires

\[
n-2\leq18\left(\sum_{e_i\neq0}3^i\right)\sqrt p
\leq9(3^k-1)\sqrt p. \qquad (8.3)
\]

\textbf{Proof.} The relation makes the associated nontrivial cubic
character product constant on every nonidentity rational point, because
the primitive orbit covers them all. Its full rational-point sum has
magnitude at least \(n-2\), while the Kummer bound and (8.2) give the
upper bound. \(\square\)

The exact integer comparisons printed in Appendix A show that
\(9(3^{78}-1)\sqrt p<n-2\), whereas the corresponding sufficient
inequality at \(k=79\) no longer holds. Thus the theorem excludes every
nonzero affine recurrence supported within \textbf{78 consecutive
positions}, indexed \(0,\ldots,77\). It does not assert existence of a
recurrence at length 79.

\section{9. Every short word occurs}\label{every-short-word-occurs}

For \(w\in\mathbf F_3^k\), let \(N_w\) count nonidentity rational
starting points for which the next \(k\) labels are \(w\). Character
orthogonality and (8.2) give

\[
\left|N_w-\frac{n-1}{3^k}\right|
\leq6(3^k-1)\sqrt p+1-3^{-k}. \qquad (9.1)
\]

\textbf{Proof.} Expand the indicator of each coordinate equality using
the three additive characters of \(\mathbf F_3\). The zero exponent
vector contributes \((n-1)/3^k\). Each nonzero vector contributes a
character sum bounded by \(s(e)\sqrt p+1\), the additional one
accounting for removing \(O\). In the sum over all exponent vectors, the
\(i\)th coordinate is nonzero in \(2\cdot3^{k-1}\) cases. Hence

\[
\frac1{3^k}\sum_e s(e)
=\frac{18}{3^k}(2\cdot3^{k-1})\sum_{i=0}^{k-1}3^i
=6(3^k-1).
\]

The \(3^k-1\) nonzero vectors contribute the remaining \(1-3^{-k}\).
\(\square\)

Each quotient starting state corresponds to six rational points with the
same word, so the number of occurrences on the quotient cycle is
\(N_w/6\). A sufficient condition for positivity is

\[
n-3^k>6\cdot3^k(3^k-1)\sqrt p. \qquad (9.2)
\]

It holds at \(k=39\) and therefore for every smaller positive length. At
\(k=40\) this estimate is insufficient. Accordingly \textbf{every
ternary word of length at most 39 occurs}. This is a consequence of the
conductor calculation, not an extrapolation from a finite simulated
prefix.

A deterministic autonomous symbol generator reproducing the sequence
would need at least \(3^{39}\) distinguishable starting states, since
two states producing different length-39 words cannot be the same state.
This is a state-count consequence, not a time lower bound: roughly 62
bits already suffice to name that many states. It does not exclude a
compact description with a large state space.

\section{10. Kummer towers and the distinction between driven and
autonomous
states}\label{kummer-towers-and-the-distinction-between-driven-and-autonomous-states}

The disjoint supports show that adjoining cube roots of
\(r_0,\ldots,r_{k-1}\) gives geometric deck group \((\mathbf Z/3)^k\).
Pullback by \(\Psi\) shifts the corresponding Kummer classes. Thus each
extra level adds genuinely new geometric branch data.

A related rational-function formulation uses \(u_j(t)=F^j(t)+28\). Its
new zero set is \(F^{-j}(-28)\), consisting of \(3^j\) distinct points.
The postcritical values of \(F\) are \(0,\infty,-7\), as follows from
\(F'/F=(t-56)/(t(t+28))\). Thus these inverse images are unramified and
the zero sets are disjoint: a shared point would make \(-28\) periodic,
whereas \(-28\mapsto0\mapsto\infty\mapsto\infty\). At a zero of \(u_j\),
a later \(u_{j+2}\) has a pole of order six, and all subsequent poles
there have that same order; they vanish modulo three. Alternatively, the
latest nonzero exponent in any finite relation is isolated at a fresh
simple zero. This proves that the classes \([u_j]\) are independent
modulo cubes. In particular, there is no fixed finite-dimensional
pullback-stable Kummer span containing \([t+28]\). This conclusion is
restricted to this generic Kummer representation; it is not a theorem
against all finite-field algorithms or all automata.

There is also a three-valued extension driven by the quotient. Let
\(D(P)=\chi_3(x(P))\). Because \(\chi_3(-\beta/3)=1\), the explicit
\(x\) map gives

\[
\frac{D(\Psi P)}{D(P)}=\chi_3(t(P)+28),\qquad
A(P_r)=\frac{D(P_{r+2})}{D(P_{r+1})}. \qquad (10.1)
\]

This is a small \textbf{fiber state driven by the quotient coordinate}.
The base coordinate still ranges over \(M\) states. It is not an
autonomous three-state generator of the ternary sequence, and does not
contradict the block bound.

\section{11. Explicit higher-genus quotients and elliptic
factors}\label{explicit-higher-genus-quotients-and-elliptic-factors}

The curve \(\mathcal C\) of Proposition 6.1 has a visible subgroup
\(A\simeq C_3\times C_3\) generated by

\[
a:(x,z,y)\mapsto(\beta x,z,y),\qquad
b:(x,z,y)\mapsto(x,\beta z,y).
\]

Negation of \(y\) gives a commuting involution. No assertion that this
subgroup is the entire automorphism group is needed.

\textbf{Proposition 11.1 (four quotients).} Put \(R=(\alpha+7)/7\). The
four order-three subgroups of \(A\) give quotient curves of genera
\(1,1,2,3\), respectively. They have the following models, interpreted
through smooth projective normalization:

\[
E:y^2=x^3+7,\qquad E_\alpha:y^2=z^3+\alpha+7,
\] \[
D_2:Y^2=7(w^3-1)(Rw^3-1),
\] \[
D_3:v^3=(y^2-7)(y^2-\alpha-7).
\]

Moreover, over \(K=\mathbf F_{p^3}\),

\[
\operatorname{Jac}(\mathcal C)\ \sim_K\
E\times E_\alpha\times\operatorname{Jac}(D_2)\times\operatorname{Jac}(D_3). \qquad (11.1)
\]

\textbf{Proof.} The quotients by \(\langle b\rangle\) and
\(\langle a\rangle\) forget \(z\) and \(x\), giving the two elliptic
curves. For \(\langle ab\rangle\), the invariant \(w=x/z\) satisfies

\[
w^3=\frac{y^2-7}{y^2-\alpha-7}.
\]

Setting \(Y=y(w^3-1)\) gives \(D_2\). Its degree-six polynomial has six
distinct roots because \(R\neq0,1\), so its genus is two. For
\(\langle ab^{-1}\rangle\), the invariant \(v=xz\) gives \(D_3\). The
degree-four polynomial in \(y\) has four distinct roots, and infinity is
also ramified in its cubic cover of \(\mathbf P^1_y\). Riemann--Hurwitz
gives \(2g-2=3(-2)+5\cdot2=4\), hence genus three. The conditions used
here follow from \(f(0)\neq0\) and \(f(-7)\neq0\).

For the isogeny, use the sum of pullback maps from the four quotient
Jacobians. On first cohomology with characteristic-zero coefficients,
every nontrivial character of \(A\) has a unique order-three kernel. Its
eigenspace therefore occurs in exactly one of the four quotient
invariant spaces. The trivial-character part is zero because
\(\mathcal C/A=\mathbf P^1_y\). These invariant spaces consequently form
a direct sum equal to \(H^1(\mathcal C)\); their dimensions are
\(2,2,4,6\). The pullback map from the product has full rank on Tate
modules, so is an isogeny. All quotient maps are defined over \(K\),
giving (11.1) over that field. The functorial pullback and norm maps on
Jacobians used here are standard {[}6{]}. \(\square\)

\textbf{Proposition 11.2 (reciprocal form and its two elliptic
quotients).} One has \(N_{K/\mathbf F_p}(R)=81\). The element \(R\) is a
square and a noncube in \(K\). Choose \(\gamma\in K\) with
\(\gamma^2=R\), and put \(S=\gamma+\gamma^{-1}\). Over the cubic
extension \(L=\mathbf F_{p^9}\), \(D_2\) is isomorphic to

\[
D_0:Y^2=7(X^6-SX^3+1). \qquad (11.2)
\]

The curve \(D_0\), itself defined over \(K\), has degree-two maps to the
elliptic curves

\[
E_+:V_+^2=7(u+2)(u^3-3u-S),
\] \[
E_-:V_-^2=7(u-2)(u^3-3u-S). \qquad (11.3)
\]

\textbf{Proof.} The norm equals \(-f(-7)/7^3=81\). Compatibility of norm
with residue characters gives \(\eta_{6,K}(R)=\eta_{6,p}(81)=\beta\):
indeed \(\chi_3(81)=\chi_3(\beta)=\beta^2\) by (1.2), and \(81\) is a
square, fixing the sextic value. Thus \(R\) is a square but not a cube.
Choose \(a_0\) with \(a_0^3=\gamma^{-1}\). It lies in \(L\), and
substituting \(w=a_0X\) gives (11.2).

The involutions \(\sigma(X,Y)=(X^{-1},Y/X^3)\) and \(-\sigma\), where
the minus is the hyperelliptic involution, have invariant functions

\[
u=X+X^{-1},\qquad V_\pm=Y(X\pm1)/X^2.
\]

Substitution gives (11.3). Both quartics are nonsingular because
\(S^2\neq4\), and each has the rational point \(u=\mp2,V_\pm=0\), so
their smooth models are elliptic curves. The degree-two quotient maps
will be denoted \(\pi_\pm\). \(\square\)

\textbf{Proposition 11.3 (a degree-three isogeny).} The elliptic curves
\(E_+\) and \(E_-\) in (11.3) are three-isogenous over \(K\). Also
\(\operatorname{Jac}(D_0)\sim_K E_+\times E_-\); consequently this
further splitting applies to \(\operatorname{Jac}(D_2)\) over \(L\).

\textbf{Proof.} Let \(J=\operatorname{Jac}(D_0)\) and let \(r\) be the
automorphism induced by \((X,Y)\mapsto(\beta X,Y)\). The quotient by
\(r\) has the conic model \(Y^2=7(Z^2-SZ+1)\), so \(1+r+r^2=0\) on
\(J\). Also \(\sigma r\sigma=r^{-1}\). The pullback and norm identities
for the two degree-two maps are

\[
\pi_+^*\pi_{+*}=1+\sigma,\quad
\pi_-^*\pi_{-*}=1-\sigma,\quad
\pi_{\pm*}\pi_\pm^*=[2].
\]

The first two eigenspaces are complementary on rational Tate modules,
giving the claimed product isogeny. Define

\[
\varphi=\pi_{-*}\,r\,\pi_+^*:E_+\longrightarrow E_-.
\]

Under the principal polarizations its dual is
\(\widehat\varphi=\pi_{+*}r^{-1}\pi_-^*\). Let \(A=\pi_{+*}r\pi_+^*\).
Since \(\sigma\) acts trivially on \(\pi_+^*E_+\) and
\(\pi_{+*}\sigma=\pi_{+*}\), conjugating by \(\sigma\) gives
\(A=\pi_{+*}r^{-1}\pi_+^*\). Therefore
\(2A=\pi_{+*}(r+r^{-1})\pi_+^*=[-2]\), and \(A=[-1]\) because an
elliptic-curve endomorphism group is torsion-free. Now

\[
\widehat\varphi\varphi
=\pi_{+*}r^{-1}(1-\sigma)r\pi_+^*
=\pi_{+*}(1-r\sigma)\pi_+^*
=[2]-A=[3].
\]

Thus \(\varphi\) is nonzero and has degree three. Every map in this
construction is defined over \(K\), and characteristic is not three.
\(\square\)

The field distinction matters. The reciprocal curve \(D_0\) and the
isogeny in Proposition 11.3 are defined over \(K\), but the change
\(D_2\simeq D_0\) uses \(a_0\in L\). The assertion that the original
factor \(\operatorname{Jac}(D_2)\) splits over \(K\) does not follow;
its twist and Frobenius descent require a separate argument. Likewise, a
full classification of the inverse-iteration tower or an unrestricted
cyclic-annihilator theorem is not claimed here.

\section{12. Conclusion}\label{conclusion}

The exact sextic target leads to a covariance equation whose nontrivial
character sums obstruct a single division polynomial and all products of
sufficiently small effective support. Along the primitive quotient
orbit, the cubic atom becomes a cyclic sequence with explicitly computed
holonomy, disjoint ramification layers, no affine recurrence in 78
consecutive positions, and occurrence of every ternary word of length at
most 39. Appendix A supplies the complete finite arithmetic underlying
these numerical thresholds.

The higher-genus constructions provide explicit quotient curves and
isogenies, not an unknown-scalar multiplication operation. Geometric
independence does not imply statistical independence on a finite orbit.
The full cyclic annihilator and the descent of the further genus-two
splitting remain separate questions. No faster private-key-recovery
algorithm follows from the statements proved here.

\section{Appendix A. Complete arithmetic
certificate}\label{appendix-a.-complete-arithmetic-certificate}

This appendix gives all data needed to certify the scalar order; it does
not rely on an external computation file. The prime coordinate modulus
and the curve parameters are the published SEC 2 parameters. The order
\(n\), the three large factors of \(n-1\), and every auxiliary prime
used below are certified by the following elementary criterion.

\textbf{Lucas criterion.} Suppose the complete factorization
\(q-1=\prod t^{e_t}\) is known, every \(t\) has already been proved
prime, and an integer \(a\) satisfies

\[
a^{q-1}\equiv1\pmod q,\qquad
\gcd(a^{(q-1)/t}-1,q)=1\quad\text{for every }t\mid q-1.
\]

Then \(q\) is prime and \(a\) has order \(q-1\) modulo \(q\). Indeed,
for every prime divisor \(\ell\) of \(q\), the displayed conditions
force the order of \(a\) modulo \(\ell\) to contain the full power of
every prime dividing \(q-1\). Thus \(q-1\mid\ell-1\), so \(\ell\geq q\)
and \(\ell=q\).

Use the parameters \(n,\lambda\) of Section 1 and \(h=\lambda+2\). The
labels in the table mean

\[
q_1=107361793816595537,\qquad q_2=174723607534414371449,
\] \[
q_3=341948486974166000522343609283189,
\] \[
r_0=29047611873442575647497758179.
\]

Starting with the prime \(2\), the rows are in increasing numerical
order. Each row states a complete factorization and a Lucas witness.
Integer multiplication, modular exponentiation, and the Euclidean
algorithm verify every row. No probabilistic primality test is required.

\begin{longtable}[]{@{}
  >{\raggedright\arraybackslash}p{(\columnwidth - 4\tabcolsep) * \real{0.3000}}
  >{\raggedright\arraybackslash}p{(\columnwidth - 4\tabcolsep) * \real{0.3000}}
  >{\centering\arraybackslash}p{(\columnwidth - 4\tabcolsep) * \real{0.4000}}@{}}
\toprule\noalign{}
\begin{minipage}[b]{\linewidth}\raggedright
\(q\)
\end{minipage} & \begin{minipage}[b]{\linewidth}\raggedright
Factorization of \(q-1\)
\end{minipage} & \begin{minipage}[b]{\linewidth}\centering
Witness \(a\)
\end{minipage} \\
\midrule\noalign{}
\endhead
\bottomrule\noalign{}
\endlastfoot
\(3\) & \(2\) & \(2\) \\
\(5\) & \(2^{2}\) & \(2\) \\
\(7\) & \(2\cdot 3\) & \(3\) \\
\(11\) & \(2\cdot 5\) & \(2\) \\
\(13\) & \(2^{2}\cdot 3\) & \(2\) \\
\(17\) & \(2^{4}\) & \(3\) \\
\(19\) & \(2\cdot 3^{2}\) & \(2\) \\
\(23\) & \(2\cdot 11\) & \(5\) \\
\(29\) & \(2^{2}\cdot 7\) & \(2\) \\
\(37\) & \(2^{2}\cdot 3^{2}\) & \(2\) \\
\(41\) & \(2^{3}\cdot 5\) & \(6\) \\
\(59\) & \(2\cdot 29\) & \(2\) \\
\(67\) & \(2\cdot 3\cdot 11\) & \(2\) \\
\(73\) & \(2^{3}\cdot 3^{2}\) & \(5\) \\
\(97\) & \(2^{5}\cdot 3\) & \(5\) \\
\(109\) & \(2^{2}\cdot 3^{3}\) & \(6\) \\
\(113\) & \(2^{4}\cdot 7\) & \(3\) \\
\(149\) & \(2^{2}\cdot 37\) & \(2\) \\
\(199\) & \(2\cdot 3^{2}\cdot 11\) & \(3\) \\
\(293\) & \(2^{2}\cdot 73\) & \(2\) \\
\(461\) & \(2^{2}\cdot 5\cdot 23\) & \(2\) \\
\(631\) & \(2\cdot 3^{2}\cdot 5\cdot 7\) & \(3\) \\
\(797\) & \(2^{2}\cdot 199\) & \(2\) \\
\(1409\) & \(2^{7}\cdot 11\) & \(3\) \\
\(1871\) & \(2\cdot 5\cdot 11\cdot 17\) & \(14\) \\
\(2011\) & \(2\cdot 3\cdot 5\cdot 67\) & \(3\) \\
\(2731\) & \(2\cdot 3\cdot 5\cdot 7\cdot 13\) & \(3\) \\
\(2861\) & \(2^{2}\cdot 5\cdot 11\cdot 13\) & \(2\) \\
\(4051\) & \(2\cdot 3^{4}\cdot 5^{2}\) & \(10\) \\
\(9349\) & \(2^{2}\cdot 3\cdot 19\cdot 41\) & \(2\) \\
\(16699\) & \(2\cdot 3\cdot 11^{2}\cdot 23\) & \(3\) \\
\(28181\) & \(2^{2}\cdot 5\cdot 1409\) & \(2\) \\
\(85831\) & \(2\cdot 3\cdot 5\cdot 2861\) & \(3\) \\
\(120233\) & \(2^{3}\cdot 7\cdot 19\cdot 113\) & \(3\) \\
\(305873\) & \(2^{4}\cdot 7\cdot 2731\) & \(3\) \\
\(1627771\) & \(2\cdot 3\cdot 5\cdot 29\cdot 1871\) & \(3\) \\
\(4681609\) & \(2^{3}\cdot 3\cdot 97\cdot 2011\) & \(23\) \\
\(44706919\) & \(2\cdot 3\cdot 797\cdot 9349\) & \(6\) \\
\(545358713\) & \(2^{3}\cdot 41\cdot 59\cdot 28181\) & \(5\) \\
\(297159362677\) & \(2^{2}\cdot 3^{2}\cdot 11\cdot 461\cdot 1627771\) &
\(2\) \\
\(q_1\) & \(2^{4}\cdot 16699\cdot 85831\cdot 4681609\) & \(3\) \\
\(q_2\) & \(2^{3}\cdot 17\cdot 59\cdot 4051\cdot 120233\cdot 44706919\)
& \(3\) \\
\(r_0\) & \(2\cdot 293\cdot 305873\cdot 545358713\cdot 297159362677\) &
\(2\) \\
\(q_3\) & \(2^{2}\cdot 3^{3}\cdot 109\cdot r_0\) & \(2\) \\
\(n\) & \(2^{6}\cdot 3\cdot 149\cdot 631\cdot q_1\cdot q_2\cdot q_3\) &
\(h\) \\
\end{longtable}

The last row uses \(h\) itself. It simultaneously proves that \(n\) is
prime and that \(h\) has order \(n-1\). Its seven proper-order tests are
precisely those for the distinct factors \(2,3,149,631,q_1,q_2,q_3\); in
particular, the certificate is not a traversal of the curve group.

The following additional checks use the same stated parameters. They can
be reproduced by the short exact-arithmetic code printed below;
exponentiation is repeated squaring. Here hexadecimal strings are
concatenated before conversion.

\begin{verbatim}
from math import gcd
p = 2**256 - 2**32 - 977
n = int('FFFFFFFFFFFFFFFFFFFFFFFFFFFFFFFE'
        'BAAEDCE6AF48A03BBFD25E8CD0364141', 16)
beta = int('7ae96a2b657c07106e64479eac3434e9'
           '9cf0497512f58995c1396c28719501ee', 16)
lam = int('5363ad4cc05c30e0a5261c028812645a'
          '122e22ea20816678df02967c1b23bd72', 16)
h, M = lam + 2, (n - 1) // 6
factors = [2, 3, 149, 631, 107361793816595537,
           174723607534414371449,
           341948486974166000522343609283189]
assert (beta*beta + beta + 1) % p == 0
assert (lam*lam + lam + 1) % n == 0
assert pow(h, n-1, n) == 1
for t in factors:
    assert gcd(pow(h, (n-1)//t, n)-1, n) == 1
assert pow(p, M, n) == 1
for t in factors:
    if M % t == 0:
        assert pow(p, M//t, n) != 1
assert pow(h, M, n) == (-lam*lam) % n
assert M % 3 == 1 and p % 9 == 7
assert pow((beta*pow(81, -1, p)) % p, (p-1)//3, p) == 1
assert pow(19+beta, (p-1)//3, p) == beta
assert pow((-beta*pow(3, -1, p)) % p, (p-1)//3, p) == 1
assert (n-2)**2 > (9*(3**78-1))**2*p
assert (n-2)**2 <= (9*(3**79-1))**2*p
assert (n-3**39)**2 > (6*3**39*(3**39-1))**2*p
assert (n-3**40)**2 <= (6*3**40*(3**40-1))**2*p
\end{verbatim}

The code uses ordinary arbitrary-precision integers; negative modular
exponents denote inverses. The four squared inequalities are used only
with positive sides. They therefore certify, without floating-point
rounding, the thresholds in the character-sum applications. These finite
arithmetic certificates support the explicitly stated parameters and
identities; they are not a substitute for the geometric proofs.

\section{Research provenance, contributions, and verification
status}\label{research-provenance-contributions-and-verification-status}

\subsection{Origin and development of the
investigation}\label{origin-and-development-of-the-investigation}

This mathematical investigation was initiated and directed by Louis
Desbiens and developed through extensive interaction with OpenAI's
ChatGPT. Louis Desbiens identified the original research questions, set
the objectives, guided the investigation, assessed its progress at a
conceptual level, and selected the material for public dissemination.

ChatGPT contributed substantially to mathematical exploration, including
alternative formulations, derivation of algebraic identities,
formulation and examination of conjectures and proposed proofs,
identification of potential counterexamples and exceptional cases,
design and execution of symbolic and numerical checks, organization of
relevant literature, and drafting and revision of the text. These
contributions went well beyond linguistic or editorial assistance.

\subsection{Attribution and scientific
authorship}\label{attribution-and-scientific-authorship}

Louis Desbiens is identified as the \textbf{human initiator and
curator}, not as an independent mathematical authority certifying the
results. OpenAI's ChatGPT is identified as the \textbf{AI mathematical
research contributor}, not as an accountable conventional scholarly
coauthor. These designations describe the provenance without implying
that the human curator personally derived or independently verified all
the arguments.

\subsection{Mathematical verification
status}\label{mathematical-verification-status}

The candidate mathematical contributions have not undergone independent
specialist validation, formal proof verification, or an exhaustive
originality and priority review. In particular, independent examination
is needed of Kummer character-sum hypotheses and conductors;
independence and branch-support arguments underlying the 78/39
thresholds; quotient genera, Jacobian splitting, isogeny and fields of
definition. Locally displayed proofs are proposed arguments. Established
results cited from public literature should be assessed separately from
new deductions; neither an AI-generated proof nor a successful
computational check is by itself an independently certified theorem.
Individual statements may be correct, partly correct, incomplete, or
incorrect.

\subsection{Responsibility and purpose of public
release}\label{responsibility-and-purpose-of-public-release}

Louis Desbiens assumes responsibility for the decision to release this
collection and for a good-faith account of its origin, contributors and
verification status. He does \textbf{not} claim to have personally
established or independently certified the correctness, completeness or
originality of the proposed mathematical results. The material is made
available for independent examination, criticism, reproducibility and
possible further development; publication does not constitute scientific
endorsement. This disclosure does not eliminate ordinary
responsibilities associated with distributing public material.

\subsection{Institutional independence, funding, and competing
interests}\label{institutional-independence-funding-and-competing-interests}

No institutional affiliation or endorsement is claimed. The
investigation was undertaken independently, without institutional
funding or resources. No external research funding was received. The
human curator reports no competing interests.

\subsection{Computational evidence and source
availability}\label{computational-evidence-and-source-availability}

Numerical evaluations helped develop the initial intuition. The public
repository contains the available scripts, archived results, and a
claim-by-claim verification coverage guide. These are scoped finite
arithmetic, symbolic, and numerical checks, not independent verification
of the general mathematical results. Historical exploratory calculations
for which independent source data are unavailable are not treated as
reproducible evidence.

\section{References}\label{references}

{[}1{]} S. D. Galbraith, \emph{Mathematics of Public Key Cryptography},
Cambridge University Press, 2012. Public author version and errata:
\url{https://www.math.auckland.ac.nz/}\textasciitilde sgal018/crypto-book/crypto-book.html.

{[}2{]} P. Deligne, ``La conjecture de Weil. II,'' \emph{Publications
Mathematiques de l'IHES} 52 (1980), 137--252. DOI: 10.1007/BF02684780.
Public text: \url{https://www.numdam.org/item/PMIHES}\_1980\_\_52\_\_137\_0/.

{[}3{]} J. S. Milne, \emph{Elliptic Curves}, second edition, World
Scientific, 2021. Public author version:
\url{https://www.jmilne.org/math/Books/EC2.pdf}.

{[}4{]} K. E. Lauter and K. E. Stange, \emph{The elliptic curve discrete
logarithm problem and equivalent hard problems for elliptic divisibility
sequences}, arXiv:0803.0728v2, 2008; published in \emph{Selected Areas
in Cryptography}, LNCS 5381, 2009, pp.~309--327. DOI:
10.1007/978-3-642-04159-4\_20. Public version:
\url{https://arxiv.org/abs/0803.0728v2}. The composition argument and the
distinction between a point function and a normalized sequence are used
explicitly in the text.

{[}5{]} Certicom Research, \emph{SEC 2: Recommended Elliptic Curve
Domain Parameters}, version 2.0, 27 January 2010, Section 2.4.1.
\url{https://www.secg.org/sec2-v2.pdf}.

{[}6{]} J. S. Milne, ``Jacobian varieties,'' in \emph{Arithmetic
Geometry}, Springer, 1986, pp.~167--212. Corrected public author
version, 12 June 2021, especially Sections 2 and 6:
\url{https://www.jmilne.org/math/xnotes/JVs.pdf}.
\end{document}
